% Meta-Diffusion: work allocation for the next stage
%
% Compile:  pdflatex work_split.tex   (twice, for the table of contents)
% Requires: amsmath amssymb booktabs tabularx enumitem xcolor hyperref caption geometry
%
% This document is a hand-out, not part of the report. build_combined.py does not read it.

\documentclass[11pt]{article}

\usepackage[a4paper,margin=1in]{geometry}
\usepackage{amsmath,amssymb}
\usepackage{booktabs}
\usepackage{tabularx}
\usepackage{array}
\usepackage{enumitem}
\usepackage{xcolor}
\usepackage[font=small,labelfont=bf]{caption}
\usepackage{hyperref}

\definecolor{shared}{HTML}{1D4E4A}    % shared structure
\definecolor{coord}{HTML}{B06A14}     % the task coordinate
\definecolor{addition}{HTML}{8C3A32}  % work that is not in the source note

\hypersetup{colorlinks=true, linkcolor=shared, urlcolor=shared, citecolor=shared}

\newcommand{\code}[1]{\texttt{\small #1}}
\newcommand{\dtask}{\Delta_{\mathrm{task}}}
\newcommand{\KT}{K_{T}}
\newcommand{\MS}{M_{S}}
\newcommand{\add}{\textcolor{addition}{\textbf{addition}}}

\newcolumntype{L}[1]{>{\raggedright\arraybackslash}p{#1}}

% Long \texttt file names carry no hyphenation points, so let TeX stretch the line
% rather than push a path into the margin.
\setlength{\emergencystretch}{3em}
\hyphenpenalty=1000
\tolerance=2000

\setlist[itemize]{leftmargin=1.4em, itemsep=2pt, topsep=3pt}
\setlist[enumerate]{leftmargin=1.6em, itemsep=2pt, topsep=3pt}

\title{\vspace{-1.5em}Meta-Diffusion: work allocation for the next stage}
\author{Three parallel packages, mapped to \emph{Immediate Next Steps} (9 September)}
\date{10 September 2026}

\begin{document}
\maketitle

\noindent
\begin{tabularx}{\textwidth}{@{}lX@{}}
\toprule
Repository & \code{CeciliaCaimi/Neural-SDE-Meta-Learning-}, branch \code{meta-diffusion-clean} \\
Baseline commit & \code{532361d} \\
Per-package brief & \code{handoff/<Name>/README.md} \\
Fixed values & \code{docs/PROTOCOL\_CARD.md} \\
\bottomrule
\end{tabularx}

\vspace{1em}

% ==========================================================================
\section{Correspondence with the source note}
% ==========================================================================

Every item below is keyed to a numbered section of \emph{Meta-Diffusion: Immediate Next
Steps}, dated 9 September. Two pieces of work are \add s that the note does not ask for;
both are marked as such here rather than folded in silently, and each is justified where
it appears.

\vspace{0.6em}
\noindent
\begin{tabularx}{\textwidth}{@{}l L{3.4cm} l X@{}}
\toprule
\textbf{Note} & \textbf{Work item} & \textbf{Owner} & \textbf{What it produces} \\
\midrule
\S1 & Generation-level task specificity (E12) & Jinjian &
  Four coordinates, one initial noise; four measurements and a paired visual grid \\
\addlinespace[2pt]
--- & Post-hoc controls (E12a), \add & Jinjian &
  The mean-coordinate substitution on the two full-width checkpoints, which no log on
  disk contains, plus the coordinate-loss profile \\
\addlinespace[2pt]
\S2 & Centred coordinate (E13, part one) & Jing Peng &
  $z_y = \bar z + \delta z_y$, the shared offset absorbed into the base predictor \\
\addlinespace[2pt]
\S3 & Four-way panel against step (E13, part two) & Jing Peng &
  $\mathcal{L}(z_y),\ \mathcal{L}(\bar z),\ \mathcal{L}(z_{\mathrm{shuffled}}),\
  \mathcal{L}(0)$; $\dtask$; basis usage and coordinate spread \\
\addlinespace[2pt]
\S3 & Calibration of $\dtask$ on stage 1, \add & Yu Cao &
  The same panel in the arm where the mechanism demonstrably works, plus $\dtask$
  against task distance \\
\addlinespace[2pt]
\S4 & Deployment sweeps & Jinjian &
  The harness, unpinned and extended; the sweep itself runs at integration \\
\addlinespace[2pt]
\S5 & Critical novelty test (E15) & Jing Peng &
  Generic latent conditioning against the reverse-dynamics basis, same transport map \\
\addlinespace[2pt]
\S6 & DomainNet gate & --- &
  Not started, and not startable until both conditions hold \\
\addlinespace[2pt]
\S7 & Excluded fixes & all three &
  A constraint, not a task. See section~\ref{sec:scope} \\
\addlinespace[2pt]
\S8 & Output for the next meeting & mapped &
  See section~\ref{sec:meeting} \\
\bottomrule
\end{tabularx}

\vspace{0.8em}

\noindent\textbf{Why the two additions.}
The note's \S1 asks for the generation test \emph{before} changing the model. E12a runs
first and costs ten minutes: the mean-coordinate substitution has never been applied to
the two 128-channel checkpoints, because the control was added to the code after every
CIFAR run had already finished. Its result may reframe \S1 before any generation code is
written.

The second addition is more consequential. The quantity \S3 makes decisive, $\dtask$, has
never been measured anywhere it is known to be positive. Stage~1 --- two-dimensional
Gaussian mixtures --- computes the correct coordinate against zero and against a shuffled
coordinate, and has never computed the mean-coordinate substitution. A zero on CIFAR
therefore has no calibration: nobody knows whether this diagnostic can produce a large
positive number at all, in this codebase, on any domain. Establishing that costs three
hours of laptop CPU.

\vspace{0.4em}
\noindent\textbf{Work completed before the split.}
Four items had to be settled in advance because all three packages depend on them: the
frozen three-corruption task family \code{artifacts/cifar100\_domainshift\_c3.json}
(the split builder previously existed for no runner to call); a \code{model\_cls} seam in
\code{training/loop.py} so the comparison arm needs no edit to the loop both arms share;
\code{--domainshift-path} on the runners; and \code{CIFAR100\_ROOT} so machines that keep
the dataset elsewhere need no code change. All are in the baseline commit.

% ==========================================================================
\section{Allocation, and why}
% ==========================================================================

\noindent
\begin{tabularx}{\textwidth}{@{}l l l X@{}}
\toprule
& \textbf{Machine} & \textbf{Branch} & \textbf{Why this machine} \\
\midrule
Jinjian & RTX 5080, local & \code{jinjian} &
  Holds the only copy of the three trained checkpoints (430\,MB each, not in git) and
  the CIFAR-100 data. Sampling 20{,}480 images needs a CUDA device, but only for minutes \\
\addlinespace[3pt]
Jing Peng & AWS GPU server & \code{jing-peng} &
  5.5 GPU-hours of training, and the only machine that can run several of them at once \\
\addlinespace[3pt]
Yu Cao & Mac, CPU only & \code{yu-cao} &
  Stage~1 is synthetic, so no dataset and no transfer; measured at 54 iterations per
  second on CPU, which is 80{,}000 steps in 24 minutes \\
\bottomrule
\end{tabularx}

\vspace{0.6em}
\noindent
The three packages are independent by construction. Everything shared was fixed before
the split; every judgement is deferred to integration; and each file in the repository
has exactly one owner. Nobody needs to wait on, or ask, anyone else.

% ==========================================================================
\section{Rules that apply to all three}
% ==========================================================================

\begin{itemize}
  \item \textbf{One owner per file.} The table is in \code{docs/PROTOCOL\_CARD.md}. If a
    file is not yours, do not edit it, even to fix an obvious bug --- record the bug in
    your results file instead.
  \item \textbf{Deliver numbers, not conclusions.} Whether the method passes is decided at
    integration, against criteria fixed before anyone sees a result.
  \item \textbf{Report absolute paired differences with 95\,\% intervals, and state the
    step count.} Do not headline \code{task\_specific\_frac}: its denominator decays
    towards zero and changes sign, so the ratio is unstable exactly where the reading
    matters.
  \item \textbf{Read nothing before the curve has levelled.} Quote the mean over the final
    fifth of training and the slope over the last 10{,}000 steps.
  \item \textbf{Pair everything.} All controls share one noising draw. Between-episode
    variance is roughly sixteen times the effect; the paired difference has a standard
    deviation sixty-five times smaller than the raw loss, so a raw loss curve cannot show
    this effect and will mislead anyone who reads it.
  \item \textbf{Results go in \code{handoff/<Name>/results/}}, which is tracked. Not
    \code{artifacts/}, which ignores \code{.txt}, \code{.png} and \code{checkpoints/}, so
    files put there vanish without an error.
\end{itemize}

% ==========================================================================
\section{Jinjian --- Measurement}
% ==========================================================================

\noindent\emph{Answers \S1; prepares \S4. About 5 hours of development, 1.5 hours of GPU.
No training.}

\begin{enumerate}
  \item \textbf{Post-hoc controls} (\add, 10 minutes of GPU). Rebuild each model from the
    checkpoint's own stored configuration, apply the moving-average weights, and evaluate
    the four coordinates on 24 validation episodes with four noise draws each, reporting
    per-episode paired differences. Then trace the coordinate-loss profile,
    $\mathcal{L}(\bar z + s\,(z_y - \bar z))$ against
    $\mathcal{L}(\bar z + s\,(z_j - \bar z))$ for $s \in \{0,\tfrac12,1,\tfrac32,2\}$.
    Flat along the correct ray and rising along wrong rays would mean the basis expresses
    distance from a centroid and nothing else.
  \item \textbf{Generation-level task specificity} (\S1). Twenty validation classes, four
    coordinates --- own, mean, shuffled, zero --- and 256 samples each, by DDIM at
    $\eta = 0$ over 50 steps, with \emph{identical initial noise across all four}. Four
    measurements: transformation consistency, semantic-class consistency, distributional
    distance, and task-identification accuracy against a chance baseline of 5\,\%. Plus a
    paired visual grid, three episodes chosen before looking at the numbers.
  \item \textbf{Sweep harness} (\S4). Unpin the checkpoint in
    \code{scripts/phase\_d\_sweeps.sh}, and add the one statistic the table does not
    currently print: the value of $\KT$ at which target-only overtakes transport.
    Coordinate space already places that crossover between 5 and 20.
\end{enumerate}

\noindent\textbf{Departure from the note, to be stated in the results file.} \S1 asks for
the kernel Inception distance. Inception is a dependency this repository does not carry,
and with 115 real images per class the estimate would be dominated by its own variance.
Section~9 of \code{ALGORITHM.md} already prescribes sliced Wasserstein or energy MMD for
this case. The standard number can be added at integration on pooled samples.

\noindent\textbf{Deliverables.} \code{E12a\_controls.txt}, \code{E12\_generation.txt},
\code{E14\_harness\_check.txt}, \code{coord\_profile.png}, \code{gen\_grid.png}.

% ==========================================================================
\section{Jing Peng --- Training}
% ==========================================================================

\noindent\emph{Answers \S2, \S3 and \S5. About 5.5 hours of development, 5.5 hours of GPU;
the runs are independent and parallelise across instances.}

\begin{enumerate}
  \item \textbf{Centred coordinate} (\S2). Write
    \[
      \hat\varepsilon_y(x,t)
      = \underbrace{\big[\hat\varepsilon_0(x,t) + B(x,t)\,\bar z\big]}_{\text{shared predictor}}
      + B(x,t)\,\delta z_y ,
    \]
    as a buffer holding a running mean, subtracted inside the single place every
    coordinate already passes through. Roughly fifteen lines; no new network. Three
    consequences must be decided rather than discovered: the null control moves from
    $z=0$ to $z=\bar z$, so \code{loss\_zero} leaves the centred panel; the bit-for-bit
    identity assertion in \code{tests/test\_score\_identity.py} is parameterised on the
    flag rather than deleted; and the drifting buffer is excluded from the
    moving-average shadow, which otherwise shadows every floating-point entry of the
    state dictionary and would leave the diagnostics reading a centre training never used.
  \item \textbf{The decisive diagnostic} (\S3). Log all four losses with per-episode paired
    differences, alongside $r_{\mathrm{basis}}$, coordinate spread and the norm of the
    running centre. Raise the diagnostic sample from 8 episodes to 24 and from one noise
    draw to four. The field names \code{delta\_task} and \code{z\_center\_norm} are an
    interface: another package's plotting script reads them.
  \item \textbf{Four runs.} The 32-channel pair first, about an hour together; then the
    128-channel pair, about two and a half hours. The centring-\emph{off} runs are not
    optional and not merely a control --- no run on disk has ever logged the
    mean-coordinate panel, and two of the four deliverables \S8 asks for are that curve.
  \item \textbf{Critical novelty test} (\S5). Generic latent conditioning entering the
    U-Net through the feature-wise modulation path that time already uses, with the same
    transport map, matched on all seven conditions the note lists. Because every
    downstream component reaches the model through one method, the substitution is a
    drop-in and the matching holds structurally rather than by bookkeeping. Report both
    parameter counts rather than asserting the budgets are equal.
\end{enumerate}

\noindent\textbf{Deliverables.} \code{E13\_panel.txt}, \code{E15\_film\_panel.txt}, and
copies of the \code{*\_log.jsonl} training logs --- easy to forget, since the originals
are written where git ignores them, and they are the primary evidence. Checkpoints are
430\,MB and stay on the server.

% ==========================================================================
\section{Yu Cao --- Positive control}
% ==========================================================================

\noindent\emph{Calibrates \S3. About 3 hours of development, 3 hours of CPU. No GPU, no
dataset, no file from anyone else's machine.}

\begin{enumerate}
  \item \textbf{The four-way panel on stage 1.} Stage~1 already reports basis usage, the
    correct coordinate against zero, a shuffled coordinate, and coordinate spread. Add the
    mean-coordinate substitution, per-task paired intervals, and the coordinate-loss
    profile. On images the profile is expected to be flat along the correct ray; here a
    clear minimum at $s=1$ is expected, and the contrast between the two is the point.
  \item \textbf{Task distance.} The stage-1 family has a dial images do not:
    \code{--family related --perturb p} interpolates from every task identical to mutually
    unrelated. The repository states a falsifiable claim --- \emph{the method needs a task
    family whose members lie farther apart than the shared backbone can absorb} --- which
    predicts that $\dtask$ rises monotonically with $p$. It has never been tested. Five
    points, two hours of CPU, and the training script already exists.
  \item \textbf{The shared plotting tool}, used by everyone at integration. Two panels:
    $\dtask$ against step with an interval band, visually dominant; below it basis usage,
    coordinate spread and the centre norm. It must tolerate every new field being absent,
    since no log currently on disk has any of them, and it must shade the final fifth of
    training, which is the only window a number may be read from.
\end{enumerate}

\noindent\textbf{Deliverables.} \code{stage1\_panel.txt}, \code{stage1\_relatedness.txt},
\code{stage1\_profile.png}, \code{delta\_task\_vs\_distance.png}.

% ==========================================================================
\section{Output for the next meeting}\label{sec:meeting}
% ==========================================================================

The four items \S8 requests, and who produces each.

\vspace{0.4em}
\noindent
\begin{tabularx}{\textwidth}{@{}l X l@{}}
\toprule
\textbf{\S8} & \textbf{Requested} & \textbf{From} \\
\midrule
1 & Generation comparison across $z_y$, $\bar z$, $z_{\mathrm{shuffled}}$, $z = 0$ & Jinjian \\
2 & One centred-coordinate run showing the four losses against training step & Jing Peng \\
3 & Basis-usage and coordinate-spread curves & Jing Peng, plotted with Yu Cao's tool \\
4 & A clear conclusion on whether task specificity survives at convergence & integration \\
\bottomrule
\end{tabularx}

\vspace{0.6em}
\noindent
Item 4 is deliberately nobody's individual deliverable. The three packages deliver
numbers; the two gates --- does generation differ across coordinates, and is
$\dtask > 0$ at convergence --- are decided afterwards, against thresholds written down
before any result was seen. Also at integration: the final deployment sweep on whichever
checkpoint passes, and Jinjian's generation instruments applied to Jing Peng's comparison
arm.

Alongside item~4, the stage-1 calibration is what makes it readable. A zero on images
means one thing if the same measurement is large and positive on Gaussian mixtures, and
something quite different if it is near zero there too.

% ==========================================================================
\section{Out of scope}\label{sec:scope}
% ==========================================================================

\S7 rules out, as the primary fix: a larger $k$; another backbone size; a more complex
transport network; confidence gates; separate source and target encoders; probabilistic
latent machinery; and healthcare- or intervention-specific components. Nothing proposed
above adds a network.

The first two are not merely discouraged --- they are already done, and negative.
Coordinate dimension $k \in \{16,32,64\}$ was swept at matched steps: the early peak rises
with $k$ and all three decay to between $0.0004$ and $0.0007$, which is the signature of a
backbone absorbing the task difference rather than of a coordinate short of capacity.
Backbone width was swept at fixed $k$: apparent gain rises 725-fold as the backbone
shrinks, and the mean-coordinate substitution then showed that none of even that was
task-specific. Both sweeps were read as successes at the time and both readings were
wrong. Repeating either would cost a day to reproduce a number already in
\code{artifacts/}.

\S6 keeps DomainNet closed until two conditions hold: $\dtask > 0$ at convergence, and
generation from $z_y$ measurably different from generation from $\bar z$. Neither has been
tested yet, and neither becomes more likely with a larger dataset.

% ==========================================================================
\section*{One open question}
% ==========================================================================

The allocation assumes the AWS instance carries a GPU, and benefits from more than one.
If it does not, Jing Peng's package has nowhere to run and the split must be revised
before anyone starts.

\end{document}
